\documentclass[12pt,a4paper]{article}
\usepackage[brazil]{babel}
\usepackage[utf8]{inputenc}
\usepackage[T1]{fontenc}
\usepackage{amsmath,amssymb}
\usepackage{geometry}
\geometry{left=3cm,right=2cm,top=3cm,bottom=2cm}
\title{Discretização da Cavidade com Tampa Deslizante}
\author{Projeto PyQt/Python}
\date{\today}

\begin{document}
\maketitle

\section*{Cavidade bidimensional com tampa deslizante}

Considere uma caixa retangular de dimensões $L_x \times L_y$, preenchida por um fluido newtoniano, incompressível e laminar. A tampa superior desliza com velocidade constante $U$, enquanto as demais paredes permanecem fixas. O escoamento é bidimensional no plano $(x,y)$.

\section*{Equações governantes}

Na forma velocidade-pressão, as equações incompressíveis são:
\begin{equation}
\frac{\partial u}{\partial x}+\frac{\partial v}{\partial y}=0.
\end{equation}
\begin{equation}
\frac{\partial u}{\partial t}+u\frac{\partial u}{\partial x}+v\frac{\partial u}{\partial y}
= -\frac{1}{\rho}\frac{\partial p}{\partial x}+\nu\left(\frac{\partial^2 u}{\partial x^2}+\frac{\partial^2 u}{\partial y^2}\right),
\end{equation}
\begin{equation}
\frac{\partial v}{\partial t}+u\frac{\partial v}{\partial x}+v\frac{\partial v}{\partial y}
= -\frac{1}{\rho}\frac{\partial p}{\partial y}+\nu\left(\frac{\partial^2 v}{\partial x^2}+\frac{\partial^2 v}{\partial y^2}\right).
\end{equation}

No projeto foi adotada a formulação vorticidade-função de corrente:
\begin{equation}
\omega = \frac{\partial v}{\partial x}-\frac{\partial u}{\partial y},
\qquad
u_x = u = \frac{\partial\psi}{\partial y},
\qquad
u_y = v = -\frac{\partial\psi}{\partial x}.
\end{equation}
Assim, resolve-se:
\begin{equation}
\nabla^2\psi = -\omega,
\end{equation}
\begin{equation}
\frac{\partial\omega}{\partial t} + u\frac{\partial\omega}{\partial x}
+ v\frac{\partial\omega}{\partial y}
= \nu\left(\frac{\partial^2\omega}{\partial x^2}+\frac{\partial^2\omega}{\partial y^2}\right).
\end{equation}

\section*{Malha node-centered}

A malha possui nós $(i,j)$ com
\begin{equation}
x_i=i\Delta x,\quad y_j=j\Delta y,
\quad \Delta x=\frac{L_x}{N_x-1},\quad \Delta y=\frac{L_y}{N_y-1}.
\end{equation}

Na abordagem \textit{node-centered}, cada incógnita é armazenada no nó. Para nós internos, o volume de controle associado ao nó $(i,j)$ possui dimensões $\Delta x\times\Delta y$. Para nós de borda, o volume é um semi-volume: por exemplo, na parede esquerda a largura é $\Delta x/2$; em cantos, o volume é um quarto de volume.

\section*{Discretização do volume interno}

Para nós internos:
\begin{equation}
\frac{\omega_{i,j}^{n+1}-\omega_{i,j}^{n}}{\Delta t}
+ u_{i,j}\left(\frac{\partial\omega}{\partial x}\right)_{i,j}
+ v_{i,j}\left(\frac{\partial\omega}{\partial y}\right)_{i,j}
= \nu\left[
\frac{\omega_{i+1,j}-2\omega_{i,j}+\omega_{i-1,j}}{\Delta x^2}
+
\frac{\omega_{i,j+1}-2\omega_{i,j}+\omega_{i,j-1}}{\Delta y^2}
\right].
\end{equation}

No código, os termos convectivos usam upwind de primeira ordem:
\begin{equation}
\left(\frac{\partial\omega}{\partial x}\right)_{i,j}=
\begin{cases}
(\omega_{i,j}-\omega_{i-1,j})/\Delta x, & u_{i,j}\geq 0,\\
(\omega_{i+1,j}-\omega_{i,j})/\Delta x, & u_{i,j}<0,
\end{cases}
\end{equation}
\begin{equation}
\left(\frac{\partial\omega}{\partial y}\right)_{i,j}=
\begin{cases}
(\omega_{i,j}-\omega_{i,j-1})/\Delta y, & v_{i,j}\geq 0,\\
(\omega_{i,j+1}-\omega_{i,j})/\Delta y, & v_{i,j}<0.
\end{cases}
\end{equation}

A equação de Poisson é discretizada por:
\begin{equation}
\frac{\psi_{i+1,j}-2\psi_{i,j}+\psi_{i-1,j}}{\Delta x^2}
+
\frac{\psi_{i,j+1}-2\psi_{i,j}+\psi_{i,j-1}}{\Delta y^2}
= -\omega_{i,j}.
\end{equation}
Isolando $\psi_{i,j}$:
\begin{equation}
\psi_{i,j}=\frac{
\Delta y^2(\psi_{i+1,j}+\psi_{i-1,j})+
\Delta x^2(\psi_{i,j+1}+\psi_{i,j-1})+
\Delta x^2\Delta y^2\omega_{i,j}
}{2(\Delta x^2+\Delta y^2)}.
\end{equation}
As velocidades nodais são recuperadas por:
\begin{equation}
u_{i,j}=\frac{\psi_{i,j+1}-\psi_{i,j-1}}{2\Delta y},
\qquad
v_{i,j}=-\frac{\psi_{i+1,j}-\psi_{i-1,j}}{2\Delta x}.
\end{equation}

\section*{Semi-volumes nas bordas e condições de contorno}

Nos nós sobre uma parede, o volume de controle possui meia espessura na direção normal à parede. A impermeabilidade impõe fluxo normal nulo. Como as paredes são linhas de corrente, adota-se:
\begin{equation}
\psi=0 \quad \text{em todas as paredes.}
\end{equation}

As condições de não escorregamento são:
\begin{align}
&\text{parede inferior: } u=0,\;v=0,\qquad
\text{parede esquerda: } u=0,\;v=0,\\
&\text{parede direita: } u=0,\;v=0,\qquad
\text{tampa superior: } u=U,\;v=0.
\end{align}

A vorticidade na parede é obtida por expansão de Taylor normal à parede, equivalente ao tratamento do semi-volume. Para uma parede horizontal parada:
\begin{equation}
\omega_w=-\frac{2\psi_P}{\Delta y^2},
\end{equation}
na qual $P$ é o primeiro nó interno adjacente à parede. Para a tampa superior móvel:
\begin{equation}
\omega_{top}=-\frac{2\psi_P}{\Delta y^2}-\frac{2U}{\Delta y}.
\end{equation}
Para paredes verticais paradas:
\begin{equation}
\omega_{left}=-\frac{2\psi_P}{\Delta x^2},
\qquad
\omega_{right}=-\frac{2\psi_P}{\Delta x^2}.
\end{equation}
Nos cantos, o código usa a média das vorticidades das duas paredes adjacentes.

\section*{Critério prático de estabilidade}

Como a integração temporal é explícita para a equação de vorticidade, recomenda-se:
\begin{equation}
\Delta t \leq \min\left( C_a\frac{h}{U},\; C_d\frac{h^2}{\nu}\right),
\quad h=\min(\Delta x,\Delta y),
\end{equation}
com $C_a<1$ e $C_d\lesssim 1/4$. O código escolhe automaticamente um passo conservador quando $\Delta t$ não é fornecido.

\end{document}
